\section{Two-option contrasts and the sharp position designs}
\label{app:general-position}
Throughout this appendix $n$ counts options and $N$ counts questions.
An order $\pi$ maps display slots to option labels. Consider the
$m=n(n-1)$ ordered pairs of distinct slots occupied by options 0 and 1.
The full-order population is uniform on these pairs. In a Latin block
each of the $n$ observed pairs is distinct, and each slot occurs once in
each coordinate. The centered weight matrix $e$ consequently vanishes
on the diagonal and has zero row and column sums. Its squared norm is
$1/n-1/m$. The space $V$ of matrices with these constraints has
dimension $d_0=m-(2n-1)$. Indeed, a linear dependence among the row and
column equations requires $a_i+b_j=0$ whenever $i\ne j$, which forces
all $a_i$ equal and all $b_j$ their negatives. For a centered matrix
$h$, its variance over orders is $\|h\|^2/m$, so its normalized risk
is $m\,\mathbb E(e^\top h)^2/\|h\|^2$, with matrices viewed as vectors.
The operator $m\mathbb E[ee^\top]$ has trace $n-2$. Every independent
extra dyad function integrated in each block removes one dimension
after projection onto $V$, proving the real-score bound of
Section~\ref{sec:theory}.

\paragraph{Why the trace bound holds for correctness.}
Write $V=V_s\oplus V_a$ for the symmetric and skew-symmetric parts.
They have dimensions $s=n(n-3)/2$ and $a=(n-1)(n-2)/2$. For symmetric
edge weights the $n$ row equations are independent since $c_i+c_j=0$
for every edge forces $c_i=0$. For skew edge weights their rank is $n-1$.
In $V_s$ assign alternating signs to both directions of the edges of a
four-cycle and zero elsewhere. In $V_a$ assign $+1$ around a triangle,
$-1$ in the opposite direction and zero elsewhere. For either witness
$h$, uniformly permute its slot labels, and consider
$F=\mathbb E[hh^\top/\|h\|^2]$. This operator is supported on the
corresponding subspace and has trace one.

On symmetric edge weights its coefficients depend only on whether two
unordered edges coincide, share a vertex, or are disjoint. At edge
$ij$ the sum on neighboring edges is $-2h_{ij}$ for every $h\in V_s$,
and the sum on disjoint edges is $h_{ij}$. Thus $F$ acts as a scalar
on $V_s$. In the skew case, exchanging the endpoints of a disjoint
edge forces its coefficient to vanish. At $ij$ the remaining
off-diagonal term is proportional to
$\sum_{t\ne i,j}(h_{it}-h_{jt})=-2h_{ij}$, again giving a scalar.
The two operators are therefore the corresponding orthogonal
projectors divided by $s$ and $a$. Selecting the two witnesses with
probabilities $s/d_0$ and $a/d_0$ gives $\mathbb E[hh^\top/\|h\|^2]
=P_V/d_0$. For every block the resulting mean normalized squared
error is $m\|e\|^2/d_0=(n-2)/d_0$. Taking the design expectation
proves Theorem~1 for correctness differences without a symmetry
assumption on the design.

\paragraph{General position-balanced budgets.}
Let $k$ be a positive multiple of $n$ with $k\le n!$, and let
$c_j$ count how often the two distinguished options occupy slot dyad
$j$ in a block. The weight errors $e_j=c_j/k-1/m$ still have zero
row and column sums. The same two ternary witnesses have conditional
mean normalized squared error $m\|e\|^2/d_0$ for each block. With
$h=k\bmod m$, the minimum of $\sum_j(c_j-k/m)^2$ over nonnegative
integer counts of sum $k$ is $h(m-h)/m$. Moving one count from a
cell to a cell smaller by at least two strictly decreases the sum,
so a minimizing vector has $h$ entries $\lceil k/m\rceil$ and all
others $\lfloor k/m\rfloor$. Applying this lower bound to every
feasible block gives the multi-budget bound in the main text.

\paragraph{Recovering the two risks from swaps.}
Randomly select the two distinguished semantic labels as well as the
slot relabeling of the witness. For two rows $x,y$, the ordered slot
dyad in $x$ is then uniform and the corresponding dyad in $y$ follows
by applying $y^{-1}x$. The inner product of a unit witness with its
image, averaged over slot relabelings, is the trace of this permutation
on $V_s$ or $V_a$ divided by the corresponding dimension, by the frame
identity above. With $x=y$ the ratio is one. With $x\ne y$ their
relative permutation is a derangement. If it has $u$ two-cycles, it
fixes $u$ unordered edges, exchanging their endpoints. Thus its trace
on all symmetric edges is $u$, and on all skew edges it is $-u$.
The symmetric vertex effects $b_i+b_j$ form the vertex permutation
representation, with trace zero on a derangement. The skew vertex
effects $b_i-b_j$ remove the constant vertex representation, giving
trace $-1$. Subtraction leaves traces $u$ and $1-u$ on $V_s$ and
$V_a$. The diagonal pairs have weight $1/n$ in the block square.
The other pairs give the two terms in Eq.~\eqref{eq:swaps}. Randomizing
the witnesses in this argument is valid even for families that are
not closed under slot relabeling.

With unrestricted distinct orders, symmetrizing a design over all
permutations of the $n!$ population indices preserves each contrast
class and yields simple random sampling. Its common risk ratio is
$(n!-n)/(n(n!-1))<1/n$. In contrast, $L_n=1/n+(n-1)/(nd_0)>1/n$.

\paragraph{Sharp upper bounds from explicit samplers.}
At each $n$ in Table~\ref{tab:template-weights}, select a template $T$
with the indicated probability. Independently shuffle its columns and
relabel its symbols uniformly, using the same resulting rows for both
models. Every displayed order has the required marginal. All template
rows are specified in \texttt{position-templates.json}. For example
the two six-row templates, written from the top down, are
\[
\begin{split}
T_1&=(012345,103254,240513,351402,425130,534021),\\
T_2&=(012345,103452,231504,345120,450231,524013).
\end{split}
\]
For templates of any size, the following calculation tests a proposed
mixture on \emph{all} real score vectors. For rows $x,y$ of $T$, two
uniform relabelings $g,h$ give the pair $(gxh,gyh)$, whose relative
permutation is $h^{-1}x^{-1}yh$. Conditional on this relative
permutation, the first row is uniform. The resulting quadratic form
under the uniform population inner product is thus a central
convolution. On the component indexed by partition $\lambda\vdash n$,
its eigenvalue before centering is
\begin{equation}
 \rho_\lambda=\frac{1}{n^2 f_\lambda}
 \sum_T w_T\sum_{x,y\in T}\chi_\lambda(x^{-1}y),
 \qquad f_\lambda=\chi_\lambda(1).
 \label{eq:character-certificate}
\end{equation}
Indeed, averaging the representation over a conjugacy class gives a
scalar matrix by Schur's lemma, and its trace determines the scalar
\citep{diaconis1981}. Centering removes the trivial component. The
largest remaining value is the real worst-case ratio. The certificate
uses integer characters, rational weights and checks all nontrivial
partitions. For each $n$ the components $(n-2,2)$ and $(n-2,1,1)$
reach the displayed maximum. Exhibiting a sampler and matching the
universal lower bound proves optimality without listing all Latin
squares of these orders.

\begin{table}[ht]
\centering
\caption{Probability of choosing each template before independent
uniform relabeling of slots and options. These are sharp upper certificates.}
\label{tab:template-weights}
\begin{tabular}{rrrr}
\toprule
$n$ & $w_1$ & $w_2$ & Largest centered eigenvalue \\
\midrule
6 & $20/133$ & $113/133$ & $4/19$ \\
7 & $25/29$ & $4/29$ & $5/29$ \\
8 & $17/41$ & $24/41$ & $6/41$ \\
9 & $37/55$ & $18/55$ & $7/55$ \\
10 & $58/71$ & $13/71$ & $8/71$ \\
\bottomrule
\end{tabular}
\end{table}

\section{Connectivity and quantitative adjacency certificates}
\label{app:connectivity}
If $X$ is a uniform row of a random block, the uniform order marginals
give $X\sim\mathrm{Unif}(S_n)$. For every profile,
\[
 v(d)=\operatorname{Var}(\mathbb E[d(X)\mid B])
      +\mathbb E\operatorname{Var}(d(X)\mid B).
\]
The first term is the risk of the sample mean before normalization.
On the centered unit sphere, equality with $v(d)$ occurs exactly when
the second term vanishes, meaning $d$ is constant on every block of
positive mass. A component indicator gives equality for every law if
the allowed graph is disconnected. If the graph is connected, uniform
sampling over the allowed blocks has full support and uniform marginals
by transitivity. Equality would force $d$ to be constant on all orders.
The maximum on the centered unit sphere is attained, so this law has
worst-case risk strictly below one. Substituting a cut indicator gives
Eq.~\eqref{eq:cut}. For the group formulation, left multiplication by
$x^{-1}$ takes any allowed block containing $x$ to a block containing
the identity. Consequently the neighbors of $x$ are $xu$, where $u$
ranges over the nonidentity rows of normalized blocks. Reachability
from $x$ is precisely its left coset of the subgroup generated by the
rows. This proves Theorem~2.

At six options the two reduced row-complete blocks are
\[
\begin{split}
B_1&=(012345,130524,204153,351402,425031,543210),\\
B_2&=(012345,135024,251403,304152,420531,543210).
\end{split}
\]
For every row $x$ here, $x(i)+x(5-i)=5$ for all zero-based slots $i$.
Under a common relabeling $g$, all rows of $gB$ therefore give the
same matching of options across opposite slots. The subgroup closure
has 24 elements. A direct enumeration at four, five and six options
generates respectively 4, 56 and 9,408 reduced Latin squares. Maximum
adjacency retains 1, 9 and 2, with subgroup orders 4, 120 and 24.

For eight, nine and ten options we use the published exhaustive
row-complete catalogues of \citet{wanlesscatalogue} (12, 12 and 492
reduced squares). Seven has no row-complete square. For a displayed
template $T$, let $q_x(i)$ be the slot of symbol $i$ in row $x$.
For a centered ternary dyad witness $h$, the normalized risk averaged
over uniform option relabelings is
\[
 U_T(h)=
 \frac{\sum_{i\ne j}\big(\sum_{x\in T}h(q_x(i),q_x(j))\big)^2}
      {n^2\sum_{u\ne v}h(u,v)^2}.
\]
The numerator enumerates the two distinguished semantic options under
the random label map. The certificate specifies six, four and six
weighted signed cycles or triangles at the respective sizes.
Direct integer evaluation gives
$\min_T\sum_h w_h U_T(h)=69/280,\ 43/243,\ 12427/99080$.
We round the final lower bound down to $1/8$ in the main table.
These statements use the completeness of the cited catalogue. For the
strict upper bounds it suffices to exhibit, respectively, two, one and
two row-complete templates whose rows generate $S_n$. In each case
the verifier exhibits a transposition and checks transitivity on
ordered pairs of distinct slots. Every transposition is then a
conjugate within the generated group. No catalogue completeness
assumption enters the upper bounds.
\section{An infinite family attaining the position bound}
\label{app:two-prime}
Let $n=2p$ with $p$ an odd prime. Put $s=n(n-3)/2$,
$a=(n-1)(n-2)/2$, $d=s+a=n^2-3n+1$, and $t=s/d<1/2$.
Choose the cyclic addition table of $\mathbb Z_n$ with probability
$w=t(n-1)/p$ and the following square otherwise. Both probabilities are positive.
In the second square, slots and symbols belong to $\mathbb F_p\times\{0,1\}$.
The $2p$ rows are the following maps, with $i\in\mathbb F_p$ and
the first coordinate calculated modulo $p$:
\[
\begin{aligned}
 A_i(x,c)&=(x+(-1)^c i,c),\\
 B_i(x,c)&=(x+(-1)^c i+c,1-c).
\end{aligned}
\]
Each row permutes the symbols. For a fixed slot, the two row families
range once over opposite fibers, proving the Latin property. Independently
permute slots and symbols uniformly after choosing the square.

\paragraph{Relative rows.}
Condition on two distinct uniformly drawn rows. Within either fiber
of row indices, the relative row translates both copies of $\mathbb F_p$
by nonzero increments, so its cycle type is $(p,p)$.
Across fibers, $A_i^{-1}B_j$ interchanges the copies and adds
$(-1)^c(i+j)+c$ to the first coordinate. Squaring adds one on both
copies, so the relative row is a single $2p$-cycle. In the cyclic
table a nonzero translation is equally likely to be any of the
$n-1$ increments. Exactly one increment yields $p$ transpositions,
and $p-1$ yield each long-cycle type. Thus in the mixture the probability
of type $2^p$ is $w/(n-1)=t/p$, the expected number of transpositions
in the cycle decomposition is $t$, and the mean sign is $-1/(n-1)$.
Indeed the probability of an even relative row in either table is
$(p-1)/(2p-1)$.

\paragraph{The low-dimensional components.}
By Eq.~\eqref{eq:character-certificate}, it suffices to bound the
off-diagonal mean character divided by degree by $1/d$ on every
nontrivial irreducible. The symmetric and skew dyad components have
dimensions $s,a$ and mean characters $t,1-t$, giving equality in
both. On a derangement with $u$ transpositions their characters are
$u,1-u$. Every row pair with $u>0$ has negative sign. The two
sign twists therefore have mean characters $-t,t-1/(n-1)$, giving
ratios at most $1/d$ since $t<1-t$. The sign character and standard
character have negative means. The sign twist of the standard character
has ratio $1/(n-1)^2<1/d$.

\paragraph{All remaining components, uniformly in the prime.}
For $n\ge256$, the low-degree classification of Rasala
\citep{rasala} gives degree $f\ge n(n-1)(n-5)/6\ge n^3/7$
after removing the trivial representation and the seven nontrivial
components treated above. On an $n$-cycle the absolute character is
at most one by the Murnaghan--Nakayama rule. On $(p,p)$ it is at
most two: in a beta-set encoding each removable $p$-hook moves one
bead to a vacant position $p$ places lower. All initially available
moves have disjoint endpoints and can be executed together. Three
moves would remove $3p$ boxes from a shape of size $2p$. After
the first move the absolute character on the remaining $p$-cycle is
at most one. On $2^p$ the Fomin--Lulov inequality \citep{fominlulov}
gives absolute character at most
\[
 \frac{2^p p!}{\sqrt{(2p)!}}\sqrt f\ \le\ (2n)^{1/4}\sqrt f .
\]
Here $\binom{2p}{p}\ge4^p/\sqrt{4p}$ follows by induction from $p=1$.
Since $t/p<1/n$, the absolute off-diagonal ratio on all these components
is at most
\[
 \frac{2}{f}+\frac{(2n)^{1/4}}{n\sqrt f}
 \le \frac{1}{n^2}\left(\frac{14}{n}+
          \frac{2^{1/4}\sqrt7}{n^{1/4}}\right)<\frac1{n^2}<\frac1d.
\]
The parenthesis is less than one already at 256 and decreases.
Adding back the diagonal gives largest centered eigenvalue
$1/n+(n-1)/(nd)=L_n$. Theorem~1 supplies the matching correctness
lower bound. This proves simultaneous optimality for every prime
$p\ge131$, without a separate search at each size.

\section{An exact attainment criterion and the twelve-option certificate}
\label{app:attainment}
For $n\ge4$, let $D=n^2-3n+1$ and let $\mathcal C$ be the set of
derangement cycle types in $S_n$. For a Latin square $T$, draw two
distinct rows $x,y$ uniformly and put $q_T(C)=\Pr(x^{-1}y\in C)$.
Let $Q_n$ be the convex hull of these histograms over Latin squares of
order $n$. Repeated histograms need only be included once.

\paragraph{Cycle-type criterion.}
The bound $L_n$ is attained simultaneously for real scores and
correctness differences if and only if there exists $q\in Q_n$ with
\begin{equation}
 \sum_{C\in\mathcal C} q(C)\chi_\lambda(C)
 \ \le\ f_\lambda/D \qquad(\lambda\vdash n,\ \lambda\ne(n)).
 \label{eq:attain-lp}
\end{equation}
These are rational linear feasibility constraints: write
$q=\sum_T w_Tq_T$, $w_T\ge0$, $\sum_Tw_T=1$.
Whenever feasible, at most $|\mathcal C|$ templates with rational
weights suffice. In particular this test uses cycle types of row
pairs rather than all $n!$ coordinates of a response. A restricted
list suffices to prove attainment; an infeasibility claim requires
control of all eligible histograms.

\emph{Proof.} Average a feasible design independently over option
and slot relabelings. This preserves the uniform marginals. Each
relabeling merely permutes the coordinates of a profile, so the
worst-case risk cannot increase in either profile class. After
averaging, the design selects a template and applies two uniform
relabelings. Conversely, every point in $Q_n$ is realized in this way.
The nonconstant eigenvalues in Eq.~\eqref{eq:character-certificate}
are $1/n+(n-1)\sum_C q(C)\chi_\lambda(C)/(nf_\lambda)$.
Comparing these values with $L_n=1/n+(n-1)/(nD)$ gives
Eq.~\eqref{eq:attain-lp} in both directions for real profiles.
The lower bound already has a correctness witness, so a feasible
solution is sharp for both classes. A rational feasible polytope
has a rational vertex. Its histogram lies in a simplex spanned by
affinely independent template histograms, at most $|\mathcal C|$
of them since their entries sum to one. Its barycentric coordinates
in that simplex are rational.

The dual form gives a certificate of impossibility. Define
$b_{T\lambda}=\sum_Cq_T(C)\chi_\lambda(C)/f_\lambda-1/D$ for
nontrivial $\lambda$. The criterion fails precisely when there is a
probability vector $u$ on these partitions such that
$\sum_\lambda u_\lambda b_{T\lambda}>0$ for every template.
Indeed the convex hull of the finite row vectors $b_T$ then misses
the closed negative orthant. Strict separation supplies a nonzero
nonnegative normal (otherwise its supremum on the orthant is infinite);
normalizing its sum gives $u$. Conversely such a $u$ excludes every
feasible mixture.

\paragraph{What balance is necessary?}
The two dyad inequalities have respective left sides $t$ and $1-t$,
where $t=\sum_Cq(C)m_2(C)$ and $m_2(C)$ counts transpositions in the
cycle decomposition. Their right sides are $s/D$ and $a/D$ with
$s+a=D$. Hence both must be equalities. This gives an initial scalar
test before the other inequalities. The same equality is necessary
even for correctness alone by Eq.~\eqref{eq:swaps}. For an exact
criterion for that class alone, replace $\chi_\lambda(C)/f_\lambda$
in the normalized inequalities by
\[
 h_z(C)=\mathbb E_{x,g}[z(x)z(xg)]/\mathbb E_x[z(x)^2]
\]
for each centered nonconstant ternary profile $z$ (center after
choosing the ternary entries), with independent uniform $x\in S_n$
and $g\in C$. Indeed the off-diagonal normalized second moment is
then $\sum_Cq(C)h_z(C)$; the diagonal contribution is unchanged.
These are again finitely many rational inequalities, though generally
far more numerous.

\paragraph{Screening the remaining degrees.}
For a candidate $q$, put $H(q)^2=\sum_C q(C)^2 n!/|C|$.
Column orthogonality gives
$\sum_\lambda(\sum_Cq(C)\chi_\lambda(C))^2=H(q)^2$.
Thus all components of degree at least $DH(q)$ satisfy the test
automatically. This bound can shorten an exact check without
assuming that the two dyad components are the worst ones.

\paragraph{Twelve options.}
The two squares in \texttt{certificates/position-templates.json}
under key \texttt{12} have total transposition counts 54 and 27 over
their 66 unordered row pairs. Draw the first square with probability
$23/109$ and the second with probability $86/109$, then shuffle
symbols and slots independently. The mean transposition count is
$(23\cdot54+86\cdot27)/(109\cdot66)=54/109=s/D$.
For this histogram $186<109H(q)<187$, so degrees at least 187 are
automatically covered. The accompanying exact certificate checks all
76 nontrivial partitions. The largest normalized risk is $10/109$
on $(10,2)$ and $(10,1,1)$; the largest remaining value is
$1873/21582<10/109$. One verifier uses the Murnaghan--Nakayama rule;
a second expands the Frobenius character formula and recomputes the
pair types from the rows. By the criterion and the correctness lower
bound, $10/109$ is the exact optimum for both classes.

